%----------------------------------------------------------------------------------------------
%-------------------------------=-     6 Matritzenzerlegung    --------------------------------
%----------------------------------------------------------------------------------------------
\section{Matritzenzerlegung}

Matritzen werden aus den folgenden Gründen zerlegt: 
\vspace{0.1cm}

\begin{itemize}
	\item Schwierige Matrix in einfachere Matritzen aufteilen
	\item Vereinfachung schwieriger Berechnungen z.B. Eigenwertzerlegung $A = T^{-1} A' T$
	\item Verienfachte Berechnung der Determinanten 
\end{itemize}
	

%----------------------------------------------------------------------------------------------
%-------------------------------------     6.1 LU-Zerlegung    --------------------------------
%----------------------------------------------------------------------------------------------
\subsection[LU-Zerlegung]{LU-Zerlegung $(A = L \cdot U)$}

\vspace{-0.2cm}
$$ \det(A) =  \det(L) \cdot \det(U)$$  

\begin{itemize}
	\item \textbf{$\bm{L}$-Matrix:} Matrix bestehnd aus Pivot-Spalten (oben rechts Nullen)
	\item \textbf{$\bm{U}$-Matrix:} Matrix mit Diagonale mit 1en, oben rechts Elemente von Gauss ohne Rückwärts-Einsetzen, 
		unten rechts Nullen 	
\end{itemize}		


\example{LU-Zerlegung}
		
$A = \begin{pmatrix}
	3  & 3  & -3 \\
	-6 & -4 & 2  \\
	3  & 13 & -18 \end{pmatrix}$ 
	\qquad \textrightarrow\ Gauss-Algorithmus durchführen
\vspace{0.3cm}
		
\begin{minipage}{0.25\columnwidth}
	\begin{tabular}{|c c c|}
		\hline
		$\cbl{3}$	& $3$  & $-3$ 	\\
		$\cbl{-6}$ 	& $-4$ & $2$ 	\\
		$\cbl{3}$	& $13$ & $-18$ 	\\
		\hline
	\end{tabular}
\end{minipage}	
\hfill
\begin{minipage}{0.23\columnwidth}
	\begin{tabular}{|c c c|}
		\hline
		$1$ & $1$ 			& $-1$  \\
		$0$ & $\cor{2}$ 	& $-4$  \\
		$0$ & $\cor{10}$ 	& $-15$ \\
		\hline
	\end{tabular}
\end{minipage}		
\hfill
\begin{minipage}{0.21\columnwidth}
	\begin{tabular}{|c c c|}
		\hline
		$1$ & $1$ & $-1$		\\
		$0$ & $1$ & $-2$ 		\\
		$0$ & $0$ & $\cgn{5}$ 	\\
		\hline
	\end{tabular}
\end{minipage}	
\hfill
\begin{minipage}{0.23\columnwidth}
	\begin{tabular}{|c c c|}
		\hline
		$\cvt{1}$	& $\cvt{1}$	& $\cvt{-1}$	\\
		$0$ 		& $\cvt{1}$ & $\cvt{-2}$ 	\\
		$0$ 		& $0$ 		& $\cvt{1}$ 	\\
		\hline
	\end{tabular}
\end{minipage}	

\vspace{0.2cm}

$$ L = \begin{pmatrix} 
	\cbl{3} & 0 & 0 \\ \cbl{-6} & \cor{2} & 0 \\ \cbl{3} & \cor{10} & \cgn{5}
\end{pmatrix}
\qquad 
U = \begin{pmatrix} 
	\cvt{1} & \cvt{1} & \cvt{-1} \\ 0 & \cvt{1} & \cvt{-2} \\ 0 & 0 & \cvt{1}  
\end{pmatrix} $$


%----------------------------------------------------------------------------------------------
%-----------------------------------     6.2 LR-Zerlegung    ----------------------------------
%----------------------------------------------------------------------------------------------
\subsection[LR-Zerlegung]{LR-Zerlegung $(A = L' \cdot R)$}

Die LR-Zerlegung (bzw. L'R-Zerlegung)  entsteht aus der LU-Zerlegung
$$ \det(A) = \det(L') \cdot \det(R) $$

\begin{tabular}{lll}
	$L' = L D^{-1} $ & $R = DU$ & $D =$ Diagonalmatrix mit Pivot-Elementen 
\end{tabular}


		
		
\example{LR-Zerlegung}

Fortsetzung des Beispiels der LU-Zerlegung!

\vspace{0.2cm}

$D = \begin{pmatrix} 
		3 & 0 & 0 \\ 0 & 2 & 0 \\ 0 & 0 & 5 
	\end{pmatrix}$ 
\quad $D^{-1} = 
	\begin{pmatrix} 
		\frac{1}{3} & 0 & 0 \\ 0 & \frac{1}{2} & 0 \\ 0 & 0 & \frac{1}{5} 
	\end{pmatrix}$
		
\vspace{0.2cm}
		
$L' = L D^{-1} =
	\begin{pmatrix} 
		3 & 0 & 0 \\ -6 & 2 & 5 \\ 3 & 10 & 5 
	\end{pmatrix}  
	\begin{pmatrix} 
		\frac{1}{3} & 0 & 0 \\ 0 & \frac{1}{2} & 0 \\ 0 & 0 & \frac{1}{5}
	\end{pmatrix} 
	= 
	\begin{pmatrix} 
		1 & 0 & 0 \\ -2 & 1 & 0 \\ 1 & 5 & 1
	\end{pmatrix}$ 
		
\vspace{0.2cm}
		
$R = DU = 
	\begin{pmatrix} 
		3 & 0 & 0 \\ 0 & 2 & 0 \\ 0 & 0 & 5 
	\end{pmatrix} 
	\begin{pmatrix} 
		1 & 1 & -1 \\ 0 & 1 & -2 \\ 0 & 0 & 1 
	\end{pmatrix}$
		

%----------------------------------------------------------------------------------------------
%--------------------------------     6.3 Cholesky-Zerlegung    -------------------------------
%----------------------------------------------------------------------------------------------
\subsection[Cholesky-Zerlegung]{Cholesky-Zerlegung $(A = L L^T)$}

\begin{itemize}
	\item Cholesky-Zerlegung existiert nur für \textbf{positiv definite, quadratische} Matritzten! 
	\item Matritzen mit negativer Determinante haben keine Cholesky-Zerlegung
	\item Matritzen mit Cholesky-Zerlegung haben Eigenwerte $\lambda_i > 0$
	\item Cholesky-Zerlegung entspricht 'Wurzel ziehen'
\end{itemize}

$$ \det(L L^T) = \det(L)^2 > 0 $$ 
		

\example{Cholesky-Zerlegung}

Die Matritzen $L$ und $L^T$ werden durch ausprobieren gefunden.

Als Vorbereitung überall wo Nullen stehen müssen, die Nullen einfüllen!\\
Anschliessend Schritt für Schritt die benötigten Werte berechnen

\vspace{0.2cm}

$A = L L^T = 
	\begin{pmatrix}
		4  & 6  & -4 \\
		6  & 10 & -7 \\
		-4 & -7 & 6 
	\end{pmatrix} 
= 
	\begin{pmatrix} 
		\cbl{2} 	& 0 		& 0 					\\
		\crd{3}		& \cor{1} 	& 0 					\\
		\cgn{-2} 	& \cvt{-1} 	& \textcolor{teal}{1}  
	\end{pmatrix} 
	\begin{pmatrix} 
		\cbl{2}	& \crd{3} 	& \cgn{-2} 				\\
		0 		& \cor{1} 	& \cvt{-1} 				\\
		0 		& 0 		& \textcolor{teal}{1} 
	\end{pmatrix} $
		
\vspace{0.2cm}
		
\textbf{Berechnung der Werte (Dokumentation Lösungsweg)}
\vspace{0.2cm}
		
\renewcommand{\arraystretch}{1.4}
\begin{tabular}{lll}
	$L_{11} = L_{11}^T$ & $? \cdot ? = ?^2 = 4$ 							& $? = \sqrt{4} = \cbl{2}$ 												\\ 
	$L_{21} = L_{12}^T$ & $\cbl{2} \cdot ? = 6$ 							& $? = 6 / \cbl{2} = \crd{3} $ 											\\
	$L_{31} = L_{13}^T$ & $\cbl{2} \cdot ? = -4$ 							& $? = -4 / \cbl{2} = \cgn{-2} $ 										\\
	$L_{22} = L_{22}^T$ & $\crd{3}^2 + ?^2 = 10$ 							& $? = \sqrt{10 - \crd{3}^2 } = \cor{1}$ 								\\	
	$L_{32} = L_{23}^T$ & $ \cgn{-2} \cdot \crd{3} + \cor{1} \cdot ? = -7 $ & $? = \frac{-7 - (\cgn{-2}) \cdot \crd{3}}{\cor{1}} = \cvt{-1} $ 		\\	
	$L_{33} = L_{33}^T$ & $\cgn{-2}^2 + \cvt{-1}^2 + ?^2 = 6 $ 				& $ ? = \sqrt{6 - (\cgn{-2})^2 - (\cvt{-1})^2} = \textcolor{teal}{1}$ 			
\end{tabular}
\renewcommand{\arraystretch}{1}


%----------------------------------------------------------------------------------------------
%-----------------------------------     6.4 QR-Zerlegung    ----------------------------------
%----------------------------------------------------------------------------------------------
\subsection[QR-Zerlegung]{QR-Zerlegung $A = Q \cdot R$}

\vspace{-0.2cm}
$$ R = Q^{-1} A = Q^T A $$  

\begin{itemize}
	\item \textbf{$\bm{Q}$-Matrix:} Orthagonale Matrix \\
		Spalten sind mit Gram-Schmidt orthagonalisierte Spalten der Matrix $A$
	\item \textbf{$\bm{R}$-Matrix:} Obere Dreiecksmatrix mit Koeffizienten der Linearkombination von Spalten von $Q$
\end{itemize}		

		
\example{QR-Zerlegung}

\begin{tabular}{lll}
	$A = 
	\begin{pmatrix}
		3 & 5 \\ 4 & 7 
	\end{pmatrix}$ 
	&
	$\vec{a}_1 = 
	\begin{pmatrix}
		3 \\ 4 
	\end{pmatrix}$ 
	&
	$\vec{a}_2 = 
	\begin{pmatrix} 
		5 \\ 7
	\end{pmatrix}$
\end{tabular}
		
\vspace{0.2cm}
		
\begin{tabular}{ll}
	$\vec{b}_1 = \frac{\vec{a1}}{\vert \vec{a}_1 \vert} = \frac{1}{5} \begin{pmatrix} 3 \\ 4 \end{pmatrix} $
	&
	$\vec{b}_2 = \frac{\vec{a}_2 - (\vec{b}_1 \bullet \vec{a}_2) \cdot \vec{b}_1}
	{\vert \vec{a}_2 - (\vec{b}_1 \bullet \vec{a}_2) \cdot \vec{b}_1 \vert} = \frac{1}{5} \begin{pmatrix} -4 \\ 3 \end{pmatrix} $ 
\end{tabular}
		
\vspace{0.2cm}

\renewcommand{\arraystretch}{1.2}
\begin{tabular}{lll}
	$Q = \frac{1}{5} \begin{pmatrix} 3 & -4 \\ 4 & 3 \end{pmatrix} $
	& &
	$R = Q^{-1} A = Q^T A = \frac{1}{5}\begin{pmatrix} 3 & 4 \\ -4 & 3 \end{pmatrix}  \begin{pmatrix} 3 & 5 \\ 4 & 7 \end{pmatrix} =
	\begin{pmatrix} 5 & \frac{43}{5} \\ 0 & \frac{1}{5}  \end{pmatrix} $
\end{tabular}
\renewcommand{\arraystretch}{1}

\pagebreak
